\documentclass[../../main.tex]{subfiles}
\begin{document}

{\bf [解答]}

% TODO: 概要図

題意の接点を$P(t,1/t)$とする．ただし$t>0$とする．このとき$(1/x)'=-1/x^2$だから
\begin{align}
  l:y=\dfrac{-1}{t^2}(x-t)+\dfrac{1}{t} \\
  \therefore l:y=\dfrac{-1}{t^2}x+\dfrac{2}{t}\equiv f(x)
\end{align}
となるので，これと$x$軸の交点は$Q(2t,0)$となる．題意の条件より
\begin{align}
  2\le 2t \quad \therefore 1 \le t \label{eq:1}
\end{align}
となる．題意の面積を$S(t)$とすると，$t$の値で場合分けして以下のようになる．

% TODO: グラフ

グラフが上図のようになるので
\begin{align}
  S(t)&=\int_1^2f(x)dx \\
  &=\left[\dfrac{-1}{2t^2}x^2+\dfrac{2}{t}x\right]_1^2  \\
  &=\dfrac{-3}{2t^2}+\dfrac{2}{t}　　\\
  &=\dfrac{-3}{2}\left(\dfrac{1}{t}-\dfrac{2}{3}\right)^2+\dfrac{2}{3}
\end{align}
となる．この時は$0<\dfrac{1}{t}\le1$だから$\dfrac{1}{t}=\dfrac{2}{3}$で最大値$\dfrac{2}{3}$をとる．
以上から求める$S(t)$の最大値は
\begin{align}
  \max S(t)=S\left(\dfrac{3}{2}\right)=\dfrac{2}{3}
\end{align}
である．

% \subparagraph{(i)$1\le2t\le2$つまり$1/2\le t\le 1$の時}

% % TODO: グラフ

% グラフが上図のようになるので
% \begin{align}
%   S(t)&=\int_1^{2t}f(x)dx  \\
%   &=\left[\dfrac{-1}{2t^2}x^2+\dfrac{2}{t}x\right]_1^{2t} \\
%   &=\dfrac{-1}{2t^2}(4t^2-1)+\dfrac{2}{t}(2t-1)  \\
%   &=2-\dfrac{2}{t}+\dfrac{1}{2t^2}　　\\
%   &=\dfrac{1}{2}(s-2)^2
% \end{align}
% となる．ただし$s=1/t$として，$1\le s\le2$である．この時の最大値は，
% $s=1$のときの$\dfrac{1}{2}$である．

% \subparagraph{(ii)$2\le2t$つまり$1\le t$の時}

\newpage
\end{document}
